\documentclass[a4paper, 11pt]{article}
\usepackage[T1]{fontenc}
\usepackage[utf8]{inputenc}
\usepackage{geometry}
\usepackage{xcolor}
\usepackage{titlesec}
\usepackage{booktabs}
\usepackage{array}
\usepackage{tabularx}
\usepackage{enumitem}
\usepackage{fancyhdr}
\usepackage{hyperref}
\usepackage[most]{tcolorbox}
\usepackage{parskip}
\usepackage{colortbl}

\geometry{top=2.5cm,bottom=2.5cm,left=2.8cm,right=2.8cm}

% Palette bianco e nero
\definecolor{primary}{HTML}{000000}
\definecolor{accent}{HTML}{333333}
\definecolor{light}{HTML}{F2F2F2}
\definecolor{success}{HTML}{000000}
\definecolor{warning}{HTML}{000000}
\definecolor{muted}{HTML}{555555}
\definecolor{boxborder}{HTML}{000000}

\titleformat{\section}{\color{primary}\large\bfseries\filright}{}{0em}{}[\color{black}\titlerule]
\titlespacing{\section}{0pt}{18pt}{8pt}
\titleformat{\subsection}{\color{accent}\normalsize\bfseries}{}{0em}{}
\titlespacing{\subsection}{0pt}{10pt}{4pt}

\pagestyle{fancy}
\fancyhf{}
\fancyhead[L]{\small\color{muted}\textbf{SCALO} -- Smart Connection \& Layover Optimizer}
\fancyhead[R]{\small\color{muted}A.A. 2025/2026}
\fancyfoot[C]{\small\color{muted}\thepage}
\renewcommand{\headrulewidth}{0.4pt}

\tcbuselibrary{breakable,skins}
\newtcolorbox{objectivebox}[2][]{enhanced,colback=light,colframe=black,coltitle=white,
  fonttitle=\bfseries\small,title={#2},arc=0pt,left=8pt,right=8pt,top=6pt,bottom=6pt,#1}
\newtcolorbox{infobox}{enhanced,colback=white,colframe=black,arc=0pt,
  left=8pt,right=8pt,top=6pt,bottom=6pt,boxrule=0.8pt}

\setlist[itemize,1]{label=\textcolor{black}{$\bullet$},leftmargin=1.4em,itemsep=3pt,parsep=0pt}

\hypersetup{colorlinks=false,linkcolor=black,urlcolor=black,
  pdftitle={SCALO -- Obiettivi Progetto Semestrale}}

\begin{document}

\begin{center}
  {\color{black}\rule{\textwidth}{2pt}}\\[10pt]
  {\fontsize{22}{26}\selectfont\bfseries\color{black}SCALO}\\[4pt]
  {\large\color{black}Smart Connection \& Layover Optimizer}\\[6pt]
  {\normalsize\color{muted}Documento degli Obiettivi -- Progetto Semestrale}\\[10pt]
  {\color{black}\rule{\textwidth}{2pt}}
\end{center}

\vspace{8pt}

\vspace{6pt}

\section{Descrizione del Progetto}

\textbf{SCALO} \`e un sito web che consente di ricercare voli tra un aeroporto di origine e uno di destinazione \emph{imponendo uno scalo desiderato} (es.\ ``voglio passare da Istanbul''). L'utente inserisce i codici IATA degli aeroporti, le date di viaggio e lo scalo obbligato; il sistema interroga le API di voli disponibili, filtra i risultati che includono tale scalo e mostra prezzi, durata totale, durata dello stopover e compagnie aeree.

Come funzionalit\`a opzionale, se la durata dello scalo supera una soglia configurabile, il sistema suggerisce attivit\`a turistiche nella citt\`a di transito tramite API dedicate (es.\ GetYourGuide).

\section{Obiettivi del Progetto}

Di seguito sono riportati gli obiettivi principali del progetto, suddivisi per priorit\`a e accompagnati dai relativi criteri di accettazione.

\begin{objectivebox}[colframe=black,colbacktitle=black]{O1 \quad|\quad MVP Funzionante}
\textbf{Descrizione:} Realizzare un prodotto minimo funzionante (MVP) che consenta la ricerca di voli andata/ritorno con scalo obbligato, restituendo risultati acquistabili tramite link partner.

\vspace{6pt}
\textbf{Criteri di accettazione:}
\begin{itemize}
  \item L'utente pu\`o inserire aeroporto di origine, destinazione e scalo obbligato tramite codice IATA.
  \item Il sistema restituisce almeno un risultato coerente con i parametri inseriti (scalo corrispondente a quello richiesto).
  \item Ogni risultato include un link verso la piattaforma di acquisto del partner API.
  \item La ricerca andata/ritorno \`e supportata con selezione delle rispettive date.
\end{itemize}
\end{objectivebox}

\vspace{8pt}

\begin{objectivebox}[colframe=black,colbacktitle=black]{O2 \quad|\quad Visualizzazione Chiara delle Informazioni}
\textbf{Descrizione:} Presentare all'utente le informazioni essenziali di ciascun volo in maniera immediata e leggibile.

\vspace{6pt}
\textbf{Criteri di accettazione:}
\begin{itemize}
  \item Ogni card di risultato mostra: prezzo totale, durata del volo end-to-end, durata dello stopover e compagnia aerea.
  \item La durata dello scalo \`e evidenziata visivamente per una consultazione rapida.
  \item I risultati sono ordinabili per prezzo e per durata totale.
  \item L'interfaccia \`e comprensibile senza istruzioni aggiuntive.
\end{itemize}
\end{objectivebox}

\vspace{8pt}

\begin{objectivebox}[colframe=black,colbacktitle=black]{O3 \quad|\quad Prestazioni e Qualit\`a UX}
\textbf{Descrizione:} Garantire un'esperienza utente fluida con tempi di risposta accettabili, anche a fronte delle latenze introdotte dalle API di terze parti.

\vspace{6pt}
\textbf{Criteri di accettazione:}
\begin{itemize}
  \item Il tempo di visualizzazione dei risultati \`e inferiore a \textbf{3 secondi} lato UI in condizioni normali.
  \item In caso di latenza elevata o errore API, l'interfaccia mostra un indicatore di caricamento e un messaggio di fallback esplicativo.
  \item Il form di ricerca \`e responsive e funzionale su dispositivi desktop e mobile.
  \item I messaggi di errore (quota API, rate limit, nessun risultato) sono chiari e informativi.
\end{itemize}
\end{objectivebox}

\vspace{8pt}

\begin{objectivebox}[colframe=black,colbacktitle=black]{O4 \quad|\quad Architettura Modulare ed Estendibile}
\textbf{Descrizione:} Progettare il codice in modo che l'aggiunta di nuovi provider API o nuove logiche di filtro richieda il minimo intervento sulle parti gi\`a esistenti.

\vspace{6pt}
\textbf{Criteri di accettazione:}
\begin{itemize}
  \item L'integrazione di un secondo provider API si esegue aggiungendo un nuovo modulo/adattatore senza modificare la logica di business centrale.
  \item Il filtro per scalo obbligato e durata minima/massima \`e separato dal layer di chiamata API.
  \item Il codice \`e documentato e strutturato in cartelle distinte per frontend, backend e configurazione.
  \item Sono presenti test su almeno 10 rotte reali (es.\ scali su AMS, IST, DOH, DXB).
\end{itemize}
\end{objectivebox}

\newpage

\section{Riepilogo Obiettivi}

\vspace{4pt}
\begin{center}
\begin{tabularx}{0.95\textwidth}{@{}
  >{\bfseries\centering\arraybackslash}p{1cm}
  >{\raggedright\arraybackslash}X
  >{\centering\arraybackslash}p{2.2cm}
  >{\centering\arraybackslash}p{2.4cm}
@{}}
\toprule
\rowcolor{light}
\textbf{ID} & \textbf{Obiettivo} \\
\midrule
O1 & MVP Funzionante                            \\[2pt]
O2 & Visualizzazione Chiara delle Informazioni  \\[2pt]
O3 & Prestazioni e Qualit\`a UX                 \\[2pt]
O4 & Architettura Modulare                      \\[2pt]
\bottomrule
\end{tabularx}
\end{center}

\section{Stack Tecnologico Previsto}

\vspace{4pt}
\begin{tabularx}{\textwidth}{@{} >{\bfseries}p{3.2cm} X @{}}
\toprule
\rowcolor{light}
\textbf{Area} & \textbf{Tecnologie} \\
\midrule
Frontend      & Vite + React (o Next.js) / Tailwind CSS \\[2pt]
Backend       & Node.js + Express (proxy API, gestione secret keys) \\[2pt]
API Voli      & Amadeus / Kiwi / Duffel / Skyscanner (da selezionare) \\[2pt]
API Attivit\`a & GetYourGuide (opzionale) \\[2pt]
Storage       & Cache in memoria per MVP; Redis (opzionale) \\[2pt]
Deploy        & Vercel / Netlify (frontend) + Render / Fly.io (backend) \\[2pt]
Compliance    & Codici IATA, rispetto ToS provider, privacy-by-design \\
\bottomrule
\end{tabularx}

\vspace{20pt}

\end{document}